\documentclass[conference]{IEEEtran}
\usepackage{times}
\usepackage[numbers]{natbib}
\usepackage{multicol}
\usepackage{graphicx}
\usepackage{booktabs}
\usepackage[bookmarks=true]{hyperref}
\usepackage{tabularx}
\usepackage{xcolor}
\usepackage{subcaption}
\newcommand{\newparagraph}[1]{{\bf #1.}}

\begin{document}

\title{Rebuttal for CoRL Submission \#416\vspace{-1cm}}
\maketitle
\thispagestyle{empty}

We thank the AC and all reviewers for constructive comments.
We clarify review questions as follows (AC=Meta Reviewer, R1=BKKY, R2=dWbG, R3=4zfY).

%基于人类行为的启发，在一些灵巧手内操作环境下，想要实现对物体的重建往往需要多相机/移动摄像头/重放置-抓取来加获取本来被遮挡的视角，因此我们做手内旋转的重建任务是试图用模拟一种更加接近人类感知世界的方式，因此我们在这个任务设置的基础上提出了ray-gpis的算法(因此controller-agnostic的不是问题)。

{\color{blue}\newparagraph{1. Additional Baselines \& Downstream Value}}

\textbf{Q1 (AC, R2, R3): How does AURORA compare with ActNeRF~\cite{actnerf}, PB-NBV~\cite{pbnbv}, and Scalable Real2Sim~\cite{pfaff2025_scalable_real2sim}?}

(1) \emph{Comparison with ActNeRF~\cite{actnerf} and PB-NBV~\cite{pbnbv}}: We compare against adapted active-planning baselines~\cite{actnerf,pbnbv} in Table~\ref{tab:main}-(a). We report F-AUC$\uparrow$ and F@5$\uparrow$ for reconstruction efficiency and final quality, {\color{red}Corr.$\uparrow$ for how well the planner ranks candidate views by their actual geometric gain}, and Time$\downarrow$ for planning efficiency. These metrics more directly isolate active-planning quality than terminal reconstruction metrics alone. Ray-GPIS ranks executable actions more effectively and plans faster by scoring reconstruction uncertainty under in-hand action constraints.
(2) \emph{Scalable Real2Sim~\cite{pfaff2025_scalable_real2sim}}: It does not involve active exploration, as all trajectories are prescribed.


\textbf{Q2 (AC, R2, R1): Is the improved mesh useful downstream?}
Yes. It is useful for not only object understanding but also in scenarios where geometry accuracy is paramount. We evaluate the benefit of AURORA using scenario in ~\cite{gualtieri2021robotic}, where we test whether
AURORA's reconstructed meshes support stable placement and subsequent regrasp.
As shown in Table~\ref{tab:main}-(b), meshes reconstructed by Ray-GPIS do so
more reliably than single-view and non-active reconstruction baselines.

% its acquisition strategy is closer to the pose-conditioned baseline in Q2.

% =========================================================
% Table 1: Planner comparison
% =========================================================
% \begin{table}[t]
% \centering
% \caption{Comparison of active view-planning algorithms.}
% \label{tab:planner}

% \setlength{\tabcolsep}{3pt}
% \begin{tabularx}{\linewidth}{Xccc}
% \toprule
% Method & Corr.$\uparrow$ & Regret$\downarrow$ & Time (s)$\downarrow$ \\
% \midrule
% \mbox{ActNeRF~\cite{actnerf}}
%   & $0.162{\pm}0.034$
%   & $0.032{\pm}0.003$
%   & $2.238{\pm}0.018$ \\
% \mbox{PB-NBV~\cite{pbnbv}}
%   & $-0.356{\pm}0.073$
%   & $0.108{\pm}0.012$
%   & $0.391{\pm}0.024$ \\
% \mbox{\textbf{Ray-GPIS}}
%   & $\mathbf{0.472{\pm}0.067}$
%   & $\mathbf{0.009{\pm}0.005}$
%   & $\mathbf{0.261{\pm}0.024}$ \\
% \bottomrule
% \end{tabularx}
% \end{table}

\textbf{Q3 (AC, R1, R2): Would a simpler pose-conditioned reorientation achieve similar performance?} Pose coverage does not imply geometric coverage. A viewpoint can be marked as visited while still yielding reconstruction holes due to depth noise or motion jitter. In this case, a pose-based approach would not revisit such regions, whereas our Ray-GPIS plans from the geometry actually acquired and can actively fill these residual gaps.


{
% may be marked as visited despite persistent hand occlusion, while depth noise, missing measurements, or motion jitter can cause observations from BundleTrack keyframes to be discarded during fusion, leaving reconstruction holes. Pose-based novelty would not revisit such regions, whereas AURORA plans from the geometry actually acquired and can actively fill these residual gaps.

}


% =========================================================
% (b) Targeted robustness ablation
% =========================================================
% \begin{tabular}{lcc}
% \toprule
% \multicolumn{3}{l}{\textbf{(b) Targeted robustness ablation}} \\
% Variant & Metric & $\Delta$ \\
% \midrule

% \multicolumn{3}{l}{\emph{Palm occlusion: Unseen R-AUC$\uparrow$}} \\
% Full Ray-GPIS
%   & [FULL] & -- \\
% w/o miss-ray interp.
%   & [FULL-0.0217] & $-0.0217$ \\

% \midrule
% \multicolumn{3}{l}{\emph{Pose jitter ($6^\circ$/3\,mm): Corr.$\uparrow$}} \\
% Full Ray-GPIS
%   & [FULL] & -- \\
% w/o RF integration
%   & [FULL-0.2304] & $-0.2304$ \\

% \bottomrule
% \end{tabular}





% \begin{table}[t]
% \centering
% \caption{Downstream success on diverse objects (30 trials/object/method).
% }
% \label{tab:downstream}
% \begin{tabularx}{\linewidth}{Xccc}
% \toprule
% Reconstruction source & Placement$\uparrow$ & Regrasp$\uparrow$ & Joint$\uparrow$ \\
% \midrule
% % Single RGB-D & 39.7 & 20.0 & 6.7 \\
% SPAR3D (oracle-align) & 30.0 & 15.0 & 5.0 \\
% TRELLIS.2 (oracle-align) & 19.7 & 7.7 & 0.0 \\
% % \midrule
% Fixed schedule & 50.0 & 44.0 & 24.0 \\
% Adapted PB-NBV & 43.3 & 40.0 & 23.0 \\
% Adapted ActNeRF & 56.0 & 72.0 & 41.7 \\
% \textbf{Full Ray-GPIS} & \textbf{57.7} & \textbf{77.7} & \textbf{45.0} \\
% \midrule
% GT mesh (oracle) & 100.0 & 99.0 & 99.0 \\
% \bottomrule
% \end{tabularx}
% \end{table}

\begin{table}[t]
\centering
\refstepcounter{table}
\label{tab:main}
\vspace{-0.08cm}
% \caption{}
\footnotesize
\setlength{\tabcolsep}{3pt}
\renewcommand{\arraystretch}{0.95}

\begin{tabularx}{\linewidth}{Xcccc}
\toprule
\multicolumn{5}{l}{\textbf{(a) Q1: Active planner comparison}} \\
Method
& F-AUC$\uparrow$
& F@5$\uparrow$
& Corr.$\uparrow$
& Time (s)$\downarrow$ \\
\midrule
ActNeRF~\cite{actnerf}
& $0.87{\pm}0.01$
& $0.97{\pm}0.01$
& $0.16{\pm}0.03$
& $2.24{\pm}0.02$ \\

PB-NBV~\cite{pbnbv}
& $0.75{\pm}0.02$
& $0.85{\pm}0.02$
& $-0.36{\pm}0.07$
& $0.39{\pm}0.02$ \\

\textbf{Ray-GPIS}
& $\mathbf{0.90{\pm}0.01}$
& $\mathbf{0.98{\pm}0.01}$
& $\mathbf{0.47{\pm}0.07}$
& $\mathbf{0.26{\pm}0.02}$ \\
\bottomrule
\end{tabularx}
\vspace{0.06cm}
% =========================================================
% (b) Downstream validation
% =========================================================
\begin{tabularx}{\linewidth}{Xccc}
\toprule
\multicolumn{4}{l}{\textbf{(b) Q2: Downstream success (\%)}} \\
Reconstruction source
  & Placement$\uparrow$
  & Regrasp$\uparrow$
  & Placement\,\&\,Regrasp$\uparrow$ \\
\midrule
SPAR3D (oracle-align)
  & 30.0 & 15.0 & 5.0 \\
TRELLIS.2 (oracle-align)
  & 19.7 & 7.7 & 0.0 \\
Fixed schedule
  & 50.0 & 44.0 & 24.0 \\
Adapted PB-NBV
  & 43.3 & 40.0 & 23.0 \\
Adapted ActNeRF
  & 56.0 & 72.0 & 41.7 \\
\textbf{Full Ray-GPIS}
  & \textbf{57.7}
  & \textbf{77.7}
  & \textbf{45.0} \\

\bottomrule
\end{tabularx}
\vspace{0.06cm}
% =========================================================
% (c) Targeted robustness ablation
% =========================================================
\begin{tabularx}{\linewidth}{Xcc}
\toprule
\multicolumn{3}{l}{\textbf{(c) Q4: Robustness of Ray-GPIS components}} \\
Method & Metric & Score \\
\midrule

\multicolumn{3}{l}{\emph{Finger occlusion: miss-ray interpolation}} \\
Full Ray-GPIS
& Unseen R-AUC$\uparrow^*$
& $\mathbf{0.5715}$ \\
w/o miss-ray interp.
& Unseen R-AUC$\uparrow^*$
& $0.5498;(-0.0217)$ \\

\midrule

\multicolumn{3}{l}{\emph{Pose errors ($6^\circ/3,\mathrm{mm}$): RF integration}} \\
Full Ray-GPIS
& Corr.$\uparrow$
& $\mathbf{0.5028}$ \\
w/o RF integration
& Corr.$\uparrow$
& $0.2724;(-0.2304)$ \\

\bottomrule
\end{tabularx}

\vspace{0.03cm}
\noindent\footnotesize{$^*$\textit{Unseen R-AUC} measures how effectively exploration recovers initially unobserved surfaces, reflecting performance under persistent occlusion.}

\vspace{-0.31cm}
\end{table}

{\color{blue}\newparagraph{2. Ablations, Robustness, Diversity \& Runtime}}

% \textbf{Q4 (R1, R3): Can the key Ray-GPIS components be ablated?}
\textbf{Q4 (R1, R3): Can the key Ray-GPIS components be ablated?}
{
Yes. We ablate two robustness components in Table~I-(c): (1) \emph{miss-ray interpolation}, which infers anchors for rays without surface hits, and (2) \emph{receptive-field (RF) integration}, which locally aggregates GP uncertainty rather than querying only the intersection point. As shown in Table~\ref{tab}-(c) and Fig.~\ref{tab}-(a), removing them degrades performance under occlusion and pose-error settings, respectively.
}



% {\color{red}Yes. We conducted an ablation in Table~\ref{tab:main}-(b). Ray-GPIS used two components to enhance robustness: (1) interpolate XXX in regions where ray does not incerption object. (2) it integrates XXXX for YYYY.  Under realistic perturbations, removing either component causes clear degradation. When ray interpolation is removed, it lead to wrong estimation in regions without ray (e.g. in palm-occulusion region), while removing integration lead to XXXX (failure pattern) pose jitter. Hence, we conclude both modules are necessary to robust uncertainty estimation.

% since each design is intended to address a specific failure mode rather than only improve nominal performance. Under realistic perturbations, removing the corresponding component causes clear degradation: miss-ray interpolation improves recovery from persistent occlusion and missing surface observations, while receptive-field integration stabilizes the uncertainty score map under pose jitter. Together, these components make Ray-GPIS robust to the imperfect observations and tracking errors encountered in real in-hand reconstruction.




\textbf{Q5 (R1, AC): How robust is AURORA to 6D pose-tracking errors?} We conducted an ablation that manually inject noise to estimated pose
on real RGB-D sequences, AURORA retains 96.3\% and 90.3\% of its
zero-error F@5 under $5^\circ$/5\,mm and $10^\circ$/10\,mm pose perturbations,
respectively.


\textbf{Q6 (R3, AC): Does the method generalize to more challenging irregular/concave objects?}
Yes. As shown in Fig.~\ref{}-(b), we additionally evaluate on the Mustard Bottle and Pitcher Base from the YCB dataset. Their F@5 scores are $0.8107{\pm}0.0695$ and $0.9274{\pm}0.0219$, respectively.

\textbf{Q7 (R3, R1, AC): What is the per-module runtime, and why replan every 6\,s?}
(1) \textbf{Interval:} The 6\,s interval is dictated by the hand controller, which needs this time to complete the tactile rotation and stabilize the grasp. (2) \textbf{Runtime:} GP update, candidate scoring, and NBV-to-action mapping take 261, 0.072, and 2.42,ms/update, respectively, while Segmentation and BundleTrack take 256.9,ms/frame.

\begin{figure}[t]
    \centering

    \fbox{\rule{0pt}{1.5cm}\rule{0.20\linewidth}{0pt}}\hfill
    \fbox{\rule{0pt}{1.5cm}\rule{0.20\linewidth}{0pt}}\hfill
    \fbox{\rule{0pt}{1.5cm}\rule{0.20\linewidth}{0pt}}\hfill
    \fbox{\rule{0pt}{1.5cm}\rule{0.20\linewidth}{0pt}}

    \vspace{0.05cm}

    \makebox[0.49\linewidth][c]{\textbf{(a)} Robustness ablation}
    \hfill
    \makebox[0.49\linewidth][c]{\textbf{(b)} New objects}

    \caption{Qualitative results for (a) robustness ablation and (b) new objects.}
    \label{fig:additional_results}
\vspace{-0.51cm}
\end{figure}
% \textbf{Q7 (R3, R1, AC): What is the per-module runtime, and why replan
% only every 6\,s?} (1) \textbf{Interval:} The 6s interval is required by the hand controller: it needs this time to respond, complete the tactile rotation, and stabilize the grasp; replanning earlier would change the command before the hand can reliably execute it. (2) \textbf{Runtime:}Segmentation and BundleTrack take 256.9,ms/frame, while GP update, candidate scoring, and NBV-to-action mapping take 261, 0.072, and 2.42,ms/update, respectively.


% \newparagraph{Other Responses}
\textbf{Q8 (R2): Why use in-hand rotation if Ray-GPIS is controller-agnostic?}
Inspired by human object inspection, we study reconstruction through in-hand reorientation rather than relying on multiple/moving cameras or place-and-regrasp to expose occluded views. Under this fixed-camera, continuously grasped setting, Ray-GPIS determines \emph{which} view to expose, while the controller determines \emph{how} to realize it.



{\small \bibliographystyle{plainnat}\bibliography{references}}

\end{document}
